\documentclass[12pt]{article}
\usepackage{amsmath}
\usepackage{amssymb}
\usepackage[letterpaper,margin=0.85in,centering]{geometry}
\usepackage{fancyhdr}
\usepackage{enumerate}
\usepackage{lastpage}
\usepackage{multicol}
\usepackage{graphicx}

\reversemarginpar

\pagestyle{fancy}
\cfoot{}
\lhead{Math 2000B}\chead{Quiz \# 12}\rhead{Thursday, 16\textsuperscript{th} April, 2015}
\rfoot{Total: 10 points}
%\chead{{\bf Name:}}
\newcommand{\points}[1]{\marginpar{\hspace{24pt}[#1]}}
\newcommand{\skipline}{\vspace{12pt}}
%\renewcommand{\headrulewidth}{0in}
\headheight 30pt

\newcommand{\di}{\displaystyle}
\newcommand{\R}{\mathbb{R}}
\newcommand{\Z}{\mathbb{Z}}
\newcommand{\aaa}{\mathbf{a}}
\newcommand{\bbb}{\mathbf{b}}
\newcommand{\ccc}{\mathbf{c}}
\newcommand{\dotp}{\boldsymbol{\cdot}}
\renewcommand{\div}[2]{#1\, |\,#2}

\begin{document}

%\author{Instructor: Sean Fitzpatrick}
\thispagestyle{fancy}
%\noindent{{\bf Name and student number:}}

\bigskip



{\bf Name: Solutions}

\bigskip

\begin{enumerate}
 \item Determine the domain and range of each of the following relations on $\R$, and sketchthe graph of each relation:

\begin{enumerate}
 \item $R=\{(x,y)\in\R\times\R\,|\, x^2+y^2=10\}$

\bigskip

Since $x^2\geq 0$ and $y^2\geq 0$, we note that $x^2=10-y^2\leq 10$, so $x^2\leq 10$, which means that $-\sqrt{10}\leq x\leq \sqrt{10}$. Similarly $-\sqrt{10}]\leq y\leq \sqrt{10}$. Thus
\begin{align*}
 \operatorname{dom}R &=\{x\in \R\,|\, -\sqrt{10}\leq x\leq \sqrt{10}\}= [-\sqrt{10},\sqrt{10}]\\
 \operatorname{range}R &=\{y\in\R\,|\, -\sqrt{10}\leq y\leq \sqrt{10}\} = [-\sqrt{10},\sqrt{10}].
\end{align*}

The graph (omitted) is that of a circle of radius $\sqrt{10}$ centred at the origin.

\item $S=\{(x,y)\in\R\times\R\,|\, x^2=y^2\}$

\bigskip

We have 
\[x^2=y^2 \leftrightarrow x^2-y^2=0\leftrightarrow (x-y)(x+y)=0 \leftrightarrow x=y \text{ or } x=-y,\]
and we see immediately that $\operatorname{dom}S=\operatorname{range}S=\R$, since $(x,x)\in S$ for all $x\in\R$.

The graph (omitted) is the pair of lines $y=x$ and $y=-x$. (It should look like a big X.)
\end{enumerate}

\item Let $A=\{1,2,3,4,5\}$. Draw a directed graph that represents the relation $R$ on $A$ defined by
\[
 R = \{(1,1),(1,2),(1,3),(2,1),(2,2),(2,3),(3,1),(3,3),(4,4),(4,5),(5,4),(5,5)\}.
\]

\bigskip


 A (somewhat poorly rendered) directed graph of the relation is given as follows:
\begin{center}
\includegraphics[width=3in]{digraph.pdf}
\end{center}

\bigskip

{\bf Bonus:} Is $R$ an equivalence relation? Why or why not?

\bigskip

It is not an equivalence relation because it is not symmetric: we have $(2,3)\in R$ but $(3,2)\notin R$.
\end{enumerate}

\end{document}